\documentclass[11pt,a4paper]{article}
\newcommand{\intestazione}{Dimero di spin --- 5. Rumore: il modello}
% =====================================================================
%  Preambolo comune ai documenti sul dimero (serie "che cosa abbiamo fatto")
%  Incluso con % =====================================================================
%  Preambolo comune ai documenti sul dimero (serie "che cosa abbiamo fatto")
%  Incluso con % =====================================================================
%  Preambolo comune ai documenti sul dimero (serie "che cosa abbiamo fatto")
%  Incluso con \input{preambolo_dimero} subito dopo \documentclass.
% =====================================================================
\usepackage[T1]{fontenc}
\usepackage[utf8]{inputenc}
\usepackage[main=italian,provide=*]{babel}
\usepackage{amsmath,amssymb,mathtools}
\usepackage{graphicx}
\usepackage{booktabs}
\usepackage{colortbl}
\usepackage{xcolor}
\usepackage{geometry}
\usepackage{placeins}
\usepackage{caption}
\usepackage{enumitem}
\usepackage{fancyhdr}
\usepackage{needspace}
\usepackage{tikz}
\usetikzlibrary{arrows.meta,positioning,shapes.geometric,calc,fit,backgrounds}
\usepackage{hyperref}
\hypersetup{colorlinks=true,linkcolor=blu,citecolor=blu,urlcolor=blu}

\geometry{margin=2.5cm}

% --- colori coerenti con quelli delle figure (palette Okabe-Ito)
\definecolor{blu}{HTML}{0072B2}
\definecolor{arancio}{HTML}{E69F00}
\definecolor{verdes}{HTML}{009E73}
\definecolor{rossos}{HTML}{D55E00}
\definecolor{violas}{HTML}{CC79A7}
\definecolor{grigetto}{HTML}{E8E8E8}

\captionsetup{font=small,labelfont=bf,width=0.92\textwidth}

\pagestyle{fancy}
\fancyhf{}
\fancyhead[L]{\small\itshape\intestazione}
\fancyhead[R]{\small\thepage}
\renewcommand{\headrulewidth}{0.3pt}

% --- notazione
\newcommand{\ket}[1]{\lvert #1\rangle}
\newcommand{\bra}[1]{\langle #1 \rvert}
\newcommand{\media}[1]{\langle #1 \rangle}
\newcommand{\vs}[1]{\mathbf{s}_{#1}}

% --- riquadro "in una riga": il messaggio da portare a casa
\newcommand{\messaggio}[1]{%
  \begin{center}
  \begin{minipage}{0.93\textwidth}
  \setlength{\fboxsep}{9pt}%
  \colorbox{grigetto}{\begin{minipage}{\dimexpr\textwidth-2\fboxsep}
  \small\textbf{In una riga.}\enspace #1
  \end{minipage}}
  \end{minipage}
  \end{center}
}

% --- riquadro "come lo sappiamo": la verifica numerica dietro un'affermazione
\newcommand{\verifica}[1]{%
  \begin{center}
  \begin{minipage}{0.93\textwidth}
  \small\textcolor{verdes}{\rule{2.2em}{1.1pt}}\enspace
  \textcolor{verdes}{\textbf{Come lo sappiamo.}}\enspace #1
  \end{minipage}
  \end{center}
  \medskip
}
 subito dopo \documentclass.
% =====================================================================
\usepackage[T1]{fontenc}
\usepackage[utf8]{inputenc}
\usepackage[main=italian,provide=*]{babel}
\usepackage{amsmath,amssymb,mathtools}
\usepackage{graphicx}
\usepackage{booktabs}
\usepackage{colortbl}
\usepackage{xcolor}
\usepackage{geometry}
\usepackage{placeins}
\usepackage{caption}
\usepackage{enumitem}
\usepackage{fancyhdr}
\usepackage{needspace}
\usepackage{tikz}
\usetikzlibrary{arrows.meta,positioning,shapes.geometric,calc,fit,backgrounds}
\usepackage{hyperref}
\hypersetup{colorlinks=true,linkcolor=blu,citecolor=blu,urlcolor=blu}

\geometry{margin=2.5cm}

% --- colori coerenti con quelli delle figure (palette Okabe-Ito)
\definecolor{blu}{HTML}{0072B2}
\definecolor{arancio}{HTML}{E69F00}
\definecolor{verdes}{HTML}{009E73}
\definecolor{rossos}{HTML}{D55E00}
\definecolor{violas}{HTML}{CC79A7}
\definecolor{grigetto}{HTML}{E8E8E8}

\captionsetup{font=small,labelfont=bf,width=0.92\textwidth}

\pagestyle{fancy}
\fancyhf{}
\fancyhead[L]{\small\itshape\intestazione}
\fancyhead[R]{\small\thepage}
\renewcommand{\headrulewidth}{0.3pt}

% --- notazione
\newcommand{\ket}[1]{\lvert #1\rangle}
\newcommand{\bra}[1]{\langle #1 \rvert}
\newcommand{\media}[1]{\langle #1 \rangle}
\newcommand{\vs}[1]{\mathbf{s}_{#1}}

% --- riquadro "in una riga": il messaggio da portare a casa
\newcommand{\messaggio}[1]{%
  \begin{center}
  \begin{minipage}{0.93\textwidth}
  \setlength{\fboxsep}{9pt}%
  \colorbox{grigetto}{\begin{minipage}{\dimexpr\textwidth-2\fboxsep}
  \small\textbf{In una riga.}\enspace #1
  \end{minipage}}
  \end{minipage}
  \end{center}
}

% --- riquadro "come lo sappiamo": la verifica numerica dietro un'affermazione
\newcommand{\verifica}[1]{%
  \begin{center}
  \begin{minipage}{0.93\textwidth}
  \small\textcolor{verdes}{\rule{2.2em}{1.1pt}}\enspace
  \textcolor{verdes}{\textbf{Come lo sappiamo.}}\enspace #1
  \end{minipage}
  \end{center}
  \medskip
}
 subito dopo \documentclass.
% =====================================================================
\usepackage[T1]{fontenc}
\usepackage[utf8]{inputenc}
\usepackage[main=italian,provide=*]{babel}
\usepackage{amsmath,amssymb,mathtools}
\usepackage{graphicx}
\usepackage{booktabs}
\usepackage{colortbl}
\usepackage{xcolor}
\usepackage{geometry}
\usepackage{placeins}
\usepackage{caption}
\usepackage{enumitem}
\usepackage{fancyhdr}
\usepackage{needspace}
\usepackage{tikz}
\usetikzlibrary{arrows.meta,positioning,shapes.geometric,calc,fit,backgrounds}
\usepackage{hyperref}
\hypersetup{colorlinks=true,linkcolor=blu,citecolor=blu,urlcolor=blu}

\geometry{margin=2.5cm}

% --- colori coerenti con quelli delle figure (palette Okabe-Ito)
\definecolor{blu}{HTML}{0072B2}
\definecolor{arancio}{HTML}{E69F00}
\definecolor{verdes}{HTML}{009E73}
\definecolor{rossos}{HTML}{D55E00}
\definecolor{violas}{HTML}{CC79A7}
\definecolor{grigetto}{HTML}{E8E8E8}

\captionsetup{font=small,labelfont=bf,width=0.92\textwidth}

\pagestyle{fancy}
\fancyhf{}
\fancyhead[L]{\small\itshape\intestazione}
\fancyhead[R]{\small\thepage}
\renewcommand{\headrulewidth}{0.3pt}

% --- notazione
\newcommand{\ket}[1]{\lvert #1\rangle}
\newcommand{\bra}[1]{\langle #1 \rvert}
\newcommand{\media}[1]{\langle #1 \rangle}
\newcommand{\vs}[1]{\mathbf{s}_{#1}}

% --- riquadro "in una riga": il messaggio da portare a casa
\newcommand{\messaggio}[1]{%
  \begin{center}
  \begin{minipage}{0.93\textwidth}
  \setlength{\fboxsep}{9pt}%
  \colorbox{grigetto}{\begin{minipage}{\dimexpr\textwidth-2\fboxsep}
  \small\textbf{In una riga.}\enspace #1
  \end{minipage}}
  \end{minipage}
  \end{center}
}

% --- riquadro "come lo sappiamo": la verifica numerica dietro un'affermazione
\newcommand{\verifica}[1]{%
  \begin{center}
  \begin{minipage}{0.93\textwidth}
  \small\textcolor{verdes}{\rule{2.2em}{1.1pt}}\enspace
  \textcolor{verdes}{\textbf{Come lo sappiamo.}}\enspace #1
  \end{minipage}
  \end{center}
  \medskip
}

\newcommand{\Tr}{\operatorname{Tr}}

\title{\bfseries Da $\ket{\psi}$ a $\rho$: mettere un rumore vero nella pipeline\\[2pt]
\large Documento 5 --- il modello di rumore}
\author{Samuele Schiavi \\ \small Tesi triennale in Fisica, Università di Parma
--- Relatore: Prof.\ Alessandro Chiesa}
\date{}

\begin{document}
\maketitle
\thispagestyle{fancy}

\section{Perché il sistema chiuso non basta più}

Tutto quello raccontato nei primi quattro documenti --- spettro esatto, VQE,
evoluzione di Trotter, correlatori --- vive su uno stato puro $\ket{\psi}$ che
evolve unitariamente. È un'idealizzazione utile per mettere a punto un metodo,
ma non è quello che succede su un dispositivo vero: un computer quantistico
reale accoppia i qubit a un ambiente che non si osserva, e ogni gate ha una
probabilità non nulla di fare qualcosa di diverso da quello che dovrebbe.

Il relatore ha chiesto di affrontare esattamente questo, come seconda parte
della tesi: prendere la stessa catena di stadi già costruita --- preparazione,
evoluzione, misura dei correlatori --- e ripeterla con un modello di rumore
realistico agganciato. Non si tocca l'Hamiltoniana: si cambia il modo in cui
il circuito viene eseguito.

\messaggio{Questa serie racconta la Parte 2 nell'ordine in cui è stata fatta,
esattamente come la serie precedente ha raccontato la Parte 1. Il filo
conduttore resta lo stesso: dallo spettro esatto (Documento 1) al VQE
(Documento 2) alla dinamica (Documento 3) ai correlatori (Documento 4) ---
qui si ripercorre la stessa catena, ma sotto rumore.}

\section{Il cambio di formalismo: perché serve $\rho$}

Una volta che il sistema è accoppiato a un ambiente che si traccia via, uno
stato puro non basta più a descriverlo: serve l'operatore densità $\rho$, e
l'evoluzione non è più un'unitaria pura ma un'operazione quantistica più
generale. In pratica, in Qiskit, questo significa passare da un simulatore a
vettore di stato a \texttt{AerSimulator(method="density\_matrix")}. Ogni
oggetto della pipeline --- energia, fedeltà, correlatore --- va riscritto in
questo linguaggio: l'energia diventa $\Tr[\rho H]$ invece di
$\bra\psi H\ket\psi$, e la fedeltà rispetto a un riferimento puro diventa
$\bra{\psi_\text{rif}}\rho\ket{\psi_\text{rif}}$.

\section{Due canali, non uno stack di rumore generico}

Il primo bivio è quanto rumore modellare. Il relatore ha confermato uno scope
preciso, non ``tutti i canali possibili'': due soli meccanismi, l'errore di
gate (canale depolarizzante, separato a 1 e 2 qubit) e l'errore di lettura
(readout). Rilassamento e defasamento ($T_1$, $T_2$) sono esclusi
esplicitamente --- una decisione dichiarata, non un'omissione tecnica.

Il canale depolarizzante a $n$ qubit ($d=2^n$) sostituisce lo stato con una
miscela fra lo stato originale e lo stato massimamente misto,
\begin{equation}
\mathcal{E}(\rho) = (1-\lambda)\,\rho + \lambda\,\frac{\mathbf 1}{d}.
\end{equation}
Le calibrazioni pubblicate non danno $\lambda$ direttamente, ma l'errore medio
di gate $\varepsilon$ misurato da randomized benchmarking; il legame fra i due
è $\varepsilon=\lambda\,(d-1)/d$, da cui $\lambda_{1q}=2\varepsilon_{1q}$ e
$\lambda_{2q}=\tfrac43\varepsilon_{2q}$.

\begin{figure}[htbp]
\centering
\includegraphics[width=0.62\textwidth]{figure/fig_depolarizzazione_bloch}
\caption{Contrazione del vettore di Bloch, $1-\lambda$, in funzione di
$\lambda$. Ai valori di calibrazione di riferimento (linee tratteggiate) la
contrazione è piccola --- meno dell'1\% --- ma non nulla: l'intero punto di
questa parte della tesi è misurare quanto questo effetto piccolo, accumulato
su molti gate, degradi i risultati a valle.}
\end{figure}

\FloatBarrier

Il readout simmetrico è descritto da un solo parametro $p$: la probabilità di
leggere il risultato sbagliato è la stessa in entrambe le direzioni. La
correzione sull'osservabile che interessa qui, $\langle Z\rangle$, è
\begin{equation}
\langle Z\rangle_\text{readout} = (1-2p)\,\langle Z\rangle_\text{ideale}
\end{equation}
--- un fattore puramente moltiplicativo. Questa singola proprietà, apparsa qui
per la prima volta, torna a essere decisiva più avanti: è la ragione per cui
il numero ottimale di passi di Trotter non dipenderà dal readout (Documento 9).

\section{Calibrazione: un chip vero, non numeri di comodo}

Piuttosto che inventare valori plausibili, i quattro numeri del modello sono
presi dalle mediane di calibrazione pubblicate per \texttt{ibm\_torino} (IBM
Heron r1, 133 qubit), da Stenger, Bazargan, Bronn, Gunlycke,
\emph{``Method for simulating open-system dynamics using mid-circuit
measurements on a quantum computer''}, arXiv:2504.15187 (2025), Appendice B:

\begin{center}
\renewcommand{\arraystretch}{1.3}
\begin{tabular}{@{}llc@{}}
\toprule
\textbf{Canale} & \textbf{Gate di riferimento} & \textbf{Valore mediano}\\
\midrule
Errore di gate, 1 qubit & \texttt{sx} & $\varepsilon_{1q}=2.9\times10^{-4}$\\
Errore di gate, 2 qubit & \texttt{CZ} & $\varepsilon_{2q}=3.8\times10^{-3}$\\
Errore di lettura & --- & $p_\text{readout}=2.3\times10^{-2}$\\
\bottomrule
\end{tabular}
\end{center}

Un punto delicato, che vale la pena rendere esplicito: il gate nativo a 2
qubit di \texttt{ibm\_torino} è il CZ, mentre il circuito di questa tesi è
scritto in una base che include il CNOT (\texttt{cx}). Usare comunque
$\varepsilon_{2q}$ del CZ come rappresentativo dell'errore a 2 qubit non è
un'approssimazione arbitraria: CNOT e CZ sono legati da un'identità
circuitale \emph{esatta}, $\text{CNOT}=(\mathbb 1\otimes H)\cdot
\text{CZ}\cdot(\mathbb 1\otimes H)$, verificata anche numericamente a
precisione macchina. L'interazione fisica a due qubit --- dove si concentra
la maggior parte dell'errore --- è la stessa identica operazione in entrambi i
casi; quello che cambia è solo un cambio di base attorno, realizzato da due
gate a 1 qubit il cui costo aggiuntivo ($\sim2\varepsilon_{1q}$, circa il
15\% di $\varepsilon_{2q}$) resta dichiarato, non nascosto. Resta comunque
corretto descrivere la scelta come ``valore rappresentativo dell'ordine di
grandezza'', non come una replica esatta dell'hardware: il randomized
benchmarking misura l'errore del CZ isolato, non del CNOT sintetizzato.

\verifica{Costruito il \texttt{NoiseModel} con questi quattro numeri
($\varepsilon_{1q}$, $\varepsilon_{2q}$, $p_\text{readout}$, e i
$\lambda$ derivati) e ricalcolato: $\lambda_{1q}=5.8000\times10^{-4}$,
$\lambda_{2q}=5.0667\times10^{-3}$, coerenti con la formula sopra a tutte le
cifre mostrate. Le istruzioni rumorose del modello risultante sono esattamente
\texttt{\{rz, sx, x, cx, measure\}}: i gate della base $\{rz,sx,x\}$ portano il
canale a 1 qubit, il \texttt{cx} quello a 2 qubit, la \texttt{measure}
l'errore di lettura --- nessun canale in più, nessun parametro nascosto.}

\section{Cosa resta fuori, dichiarato}

Per onestà di scope, come già nella serie precedente: $T_1/T_2$ sono esclusi
per decisione del relatore; ogni canale è applicato indipendentemente qubit
per qubit (nessun crosstalk); il readout è simmetrico in questo modello di
base (l'asimmetria è un'estensione a sé, Documento 11); la calibrazione è un
solo punto nel tempo, non un modello che tenga conto di quando viene eseguito
il circuito.

\messaggio{Il modello è deliberatamente minimale: due canali, quattro numeri,
calibrati su un chip reale. Questa scelta ha una conseguenza diretta su tutto
quello che segue --- ``rumoroso'', in questa tesi, significa esattamente
questi due canali a questi valori (o a quelli esplicitamente scansionati nel
Documento 9), non una replica generica di ``un computer quantistico
rumoroso''.}

\end{document}
